\section{First-Order Differential Equations}
\label{sec:ch4-first-order}

A first-order ODE relates a function to its first derivative. Although simple in form, such equations already model decay, charging, cooling, relaxation, transport, and population growth.

\subsection{Separable equations}
An equation is separable when it can be written
\[
\frac{dy}{dt}=g(t)h(y).
\]
Where $h(y)\neq0$, rearrange and integrate:
\[
\int \frac{dy}{h(y)}=\int g(t)\,dt+C.
\]
The method is justified by the chain rule, not by treating $dy$ and $dt$ as independent numbers.

\begin{workedexample}
A sample contains $N_0$ unstable nuclei and obeys $\dot N=-\lambda N$ with $\lambda>0$. Find $N(t)$ and the half-life.
\begin{solution}
Separate variables:
\[
\frac{dN}{N}=-\lambda\,dt.
\]
Integrating gives $\ln N=-\lambda t+C$, hence
\[
N(t)=N_0e^{-\lambda t}.
\]
The half-life $t_{1/2}$ satisfies $N(t_{1/2})=N_0/2$, so
\[
\boxed{t_{1/2}=\frac{\ln2}{\lambda}}.
\]
\end{solution}
The inverse rate $\tau=1/\lambda$ is the characteristic decay time.
\end{workedexample}

\subsection{First-order linear equations}
The standard form is
\[
\dot y+p(t)y=q(t).
\]
Multiplication by the integrating factor
\[
\mu(t)=\exp\left(\int p(t)\,dt\right)
\]
turns the left side into a product derivative:
\[
\frac{d}{dt}\bigl(\mu y\bigr)=\mu q.
\]
Therefore
\[
\boxed{y(t)=\frac{1}{\mu(t)}\left[C+\int^t\mu(s)q(s)\,ds\right]}.
\]

\begin{workedexample}
A body at temperature $T(t)$ is placed in an environment at constant temperature $T_a$. Newton's cooling law is
\[
\dot T=-k(T-T_a),\qquad k>0.
\]
For $T(0)=T_0$, determine $T(t)$.
\begin{solution}
Set $u=T-T_a$. Then $\dot u=-ku$, so
\[
u(t)=(T_0-T_a)e^{-kt}.
\]
Hence
\[
\boxed{T(t)=T_a+(T_0-T_a)e^{-kt}}.
\]
\end{solution}
The difference from equilibrium decays exponentially; the ambient temperature is a stable fixed point.
\end{workedexample}

\subsection{Equilibria and qualitative behavior}
For an autonomous equation $\dot y=f(y)$, an equilibrium $y_*$ satisfies $f(y_*)=0$. The sign of $f$ on either side determines whether neighboring solutions move toward or away from $y_*$. This one-dimensional phase-line analysis anticipates the stability theory of systems.

For the logistic equation
\[
\dot P=rP\left(1-\frac{P}{K}\right),
\]
the equilibria are $P=0$ and $P=K$. For $r>0$, small positive populations grow away from zero and toward the carrying capacity $K$.

\begin{workedexample}
Solve the logistic initial-value problem $\dot P=rP(1-P/K)$, $P(0)=P_0>0$.
\begin{solution}
Partial fractions give
\[
\frac{1}{P(1-P/K)}=\frac{1}{P}+\frac{1}{K-P}.
\]
Integration yields
\[
\ln\frac{P}{K-P}=rt+C.
\]
Using the initial condition,
\[
\boxed{P(t)=\frac{K}{1+\left(\frac{K-P_0}{P_0}\right)e^{-rt}}}.
\]
\end{solution}
The early-time behavior is approximately exponential, while nonlinear feedback produces saturation.
\end{workedexample}

\subsection{Direction fields}
The differential equation $\dot y=f(t,y)$ assigns a slope to every point in the $(t,y)$ plane. Drawing short line segments of slope $f(t,y)$ produces a direction field. A solution curve is tangent to this field everywhere.

\begin{figure}[htbp]
\centering
\begin{tikzpicture}[x=1.0cm,y=.65cm]
\draw[->] (-3.4,0)--(3.6,0) node[right] {$t$};
\draw[->] (0,-3.4)--(0,3.5) node[above] {$y$};
\foreach \x in {-3,-2.5,...,3}{\foreach \y in {-3,-2.5,...,3}{\pgfmathsetmacro{\ang}{atan(-\y)}\draw[gray,rotate around={\ang:(\x,\y)}] (\x-.16,\y)--(\x+.16,\y);}}
\draw[chapterblue,thick,domain=-3:3,samples=80] plot (\x,{2*exp(-\x-3)});
\draw[chapterblue,thick,domain=-3:3,samples=80] plot (\x,{-2*exp(-\x-3)});
\draw[chapterblue,thick] (-3,0)--(3,0);
\end{tikzpicture}

\caption{A direction field for $\dot y=-y$ and representative solution curves. Local slopes organize the global family.}
\label{fig:ch4-slope-field}
\end{figure}

\begin{physicalbox}[Relaxation as universal behavior]
Many systems near stable equilibrium obey $\dot u=-u/\tau$: a capacitor voltage, a temperature difference, a chemical imbalance, or a small perturbation of a damped device. The microscopic mechanisms differ, but the same linear relaxation equation produces the same exponential law.
\end{physicalbox}
